\section{The drift restriction}\label{sec:drift}

\begin{assumption}[Martingale property]\label{ass:mart}
Under $\Q$ every gridded instrument price is a martingale in $t$. Since the
gridded payoffs $\pi^i$ span $\R^{N+2}$, this is equivalent to each density
vector being a martingale,
\begin{equation}\label{eq:qdyn}
  \dd q^j_t=\psi^j_t\,\dd W_t,\qquad \psi^j_t\in\R^{(N+2)\times d}.
\end{equation}
\end{assumption}

This is the only assumption behind the drift restriction, and it carries no
regularity. Each
$q^j$ is a fixed linear combination of traded butterflies, so \eqref{eq:qdyn} is
the martingale representation theorem applied to traded prices, and nothing has
to commute with a stochastic integral. This is the substantive advantage of a
discrete code-book over its continuum analogue, where the corresponding step
requires differentiating a stochastic field in strike and maturity and must be
licensed by standing regularity conditions. Theorem~\ref{thm:pointwise} below
needs only $q^j>0$.

\begin{theorem}[Drift restriction]\label{thm:drift}
Let $\dd\dlv^j_t=\alpha^j_t\dd t+\beta^j_t\dd W_t$. Under
Assumption~\ref{ass:mart}, for each $j=1,\dots,M$,
\begin{align}
  \text{(volatility)}\qquad
  &A_j(\dlv^j)\,\psi^j\;=\;\psi^{j-1}-G_j(q^j)\,\beta^j,
  \label{eq:volsolve}\\[4pt]
  \text{(drift)}\qquad
  &G_j(q^j)\,\alpha^j\;=\;-\sum_{f=1}^{d}G_j\bigl(\psi^{j,f}\bigr)\beta^{j,f},
  \label{eq:driftsolve}
\end{align}
marched forward in $j$ from $\psi^0$, the volatility of the initial density,
which is determined by the spot dynamics.
\end{theorem}

\begin{proof}
Apply It\^o to $A_j(\dlv^j_t)q^j_t=q^{j-1}_t$, using \eqref{eq:affine} to write
the left side as $q^j+G_j(q^j)\dlv^j$ and noting that $G_j$ is bilinear. The
$\dd W$ terms give $\psi^j+G_j(\psi^j)\dlv^j+G_j(q^j)\beta^j=\psi^{j-1}$, which
is \eqref{eq:volsolve} by \eqref{eq:affine}. The $\dd t$ terms give
$G_j(q^j)\alpha^j+\sum_f G_j(\psi^{j,f})\beta^{j,f}=0$, since $q^{j-1}$ and
$q^{j}$ are driftless by Assumption~\ref{ass:mart} and the only quadratic
covariation is between $G_j(\cdot)$ and $\dlv^j$, both driven by $W$.
\end{proof}

Equation \eqref{eq:volsolve} is not ours. Compared with the pricing recursion
$A_j(\dlv^j)q^j=q^{j-1}$ it is \emph{the same operator with one added source},
the structure of Carmona and Nadtochiy \cite{CarmonaNadtochiy2008}. The drift
restriction below is theirs in the continuum, and Wissel's \cite{Wissel2007} on
this lattice; we claim no priority for either.

What the fixed lattice changes is threefold. Their restriction carries a term in
the derivative of the code-book along time to maturity; a lattice of fixed expiry
dates has no such derivative, and here that term is absent. The relation is algebraic rather than
differential, so neither the high-order regularity their statement requires nor
the singular behaviour as time to maturity vanishes that it is built to
accommodate can arise here, the front slab having positive width. And the
discrete system is over-determined, $N+2$ conditions on $N$ unknowns, which is
what the counting of Remark~\ref{cor:counting} squares; that counting is what
turns a differential relation into an identity, and it has no continuum
counterpart. Appendix~\ref{app:priority} carries both identifications in the
original notation.

\begin{theorem}[The drift is a pointwise covariance]\label{thm:pointwise}
Define the relative density volatility
\begin{equation}\label{eq:eta}
  \eta^{j,i}\;:=\;\frac{\psi^{j,i}}{q^{j,i}},
  \qquad i=0,\dots,N+1,
\end{equation}
that is, the relative volatility of $q^{j,i}$. Only the interior components enter
\eqref{eq:pointwise}, since $\dlv$ lives on the interior nodes; the boundary
components are needed in Section~\ref{sec:exist}. Then \eqref{eq:driftsolve} is solvable
for every $\beta$ and every state, its solution is unique, and it is given
componentwise by
\begin{equation}\label{eq:pointwise}
  \alpha^{j,i}\;=\;-\ip{\eta^{j,i}}{\beta^{j,i}}
  \qquad\Longleftrightarrow\qquad
  \dd\lambda^{j,i}=\Bigl[-\ip{\eta^{j,i}}{b^{j,i}}-\tfrac12\abs{b^{j,i}}^{2}\Bigr]\dd t
   +b^{j,i}\,\dd W_t,
\end{equation}
where $\lambda=\log\dlv$ and $b=\beta/\dlv$.
\end{theorem}

\begin{proof}
By Lemma~\ref{lem:factor} both sides of \eqref{eq:driftsolve} factor through the
same fixed matrix:
\[
  \widehat G_j\,\mathrm{diag}(q^j)\,\alpha^j
  \;=\;-\sum_f\widehat G_j\,\mathrm{diag}(\psi^{j,f})\,\beta^{j,f} .
\]
Since $\widehat G_j$ is injective it cancels, leaving
$\mathrm{diag}(q^j)\alpha^j=-\sum_f\mathrm{diag}(\psi^{j,f})\beta^{j,f}$, which
is \eqref{eq:pointwise} after dividing by $q^{j,i}>0$. Existence and uniqueness
follow from the same cancellation. The logarithmic form is It\^o applied to
$\lambda=\log\dlv$, using $\alpha/\dlv=-\ip{\eta}{b}$.
\end{proof}

Two features of \eqref{eq:pointwise} are worth drawing out. The first is its
shape. It is the exact analogue of the HJM restriction $\alpha=\sigma\int\sigma$:
the drift of the log discrete local variance is minus the covariance of its own
volatility with the volatility of the log density, plus an It\^o convexity term.
Where HJM's derived object is an integral of the volatility across maturities,
ours is $\eta$, obtained from $b$ by the single forward recursion
\eqref{eq:volsolve}. The second is the cost. Once $\eta$ is known no matrix is inverted: no Jacobian
of the pricing map is ever formed and no rank condition on it appears. And since
$\psi$ is linear in $\beta$ and $\alpha$ is bilinear in $(\beta,\psi)$, the drift
is an exact quadratic form in $\beta$, as in HJM; Theorem~\ref{thm:wing} is what
rules out the corresponding finite-time explosion.

The driftless point is the classical model: if $\beta\equiv0$ and $\psi^0=0$ then
$\psi\equiv0$ and $\alpha\equiv0$, so the discrete local variance surface is
frozen, and with $\psi^0\ne0$ the densities still move, driven by the spot alone,
with $\alpha\equiv0$ persisting.

\begin{proposition}[Spot consistency]\label{prop:spot}
The spot is $S_t=K^{\!\top}q^{j}_t$, the same for every $j$ by
Lemma~\ref{lem:null}. Since $K^{\!\top}A_j=K^{\!\top}$ and
$K^{\!\top}\widehat G_j=0$, equation \eqref{eq:volsolve} gives
\[
  \one^{\!\top}\psi^{j}=\one^{\!\top}\psi^{0}=0
  \qquad\text{and}\qquad
  K^{\!\top}\psi^{j}=K^{\!\top}\psi^{0}=\sigma^{S}_tS_t
  \quad\text{for every }j.
\]
So the spot's volatility is the single scalar $K^{\!\top}\psi^{0}$, it is
automatically the same at every expiry, and no further condition links the spot
to the surface. Taking $\psi^{0}=0$ freezes the spot and recovers
$K^{\!\top}\psi^{j}=0$.
\end{proposition}

Proposition~\ref{prop:spot} is the discrete form of Carmona and Nadtochiy's spot
consistency condition: the spot and its volatility are both \emph{read off} the
surface rather than modelled, $S=K^{\!\top}q^j$ and $\sigma^SS=K^{\!\top}\psi^j$
for any $j$.

The drift restriction is used below through the following chain, which is what a
model actually evaluates:
\begin{equation}\label{eq:pipeline}
  b\;\xrightarrow{\ \eqref{eq:volsolve}\ }\;\psi^1,\dots,\psi^M
  \;\xrightarrow{\ \eta=\psi/q\ }\;\alpha ,
\end{equation}
a forward march of $d$ tridiagonal solves per expiry, $O(NMd)$ in total. The
volatility structure is the only input; the drift is computed, never fitted.
